\chapter{Tip I--Tip II Hata, Test Gücü ve Tek Örneklem t Testi}
\label{ch:power-one-sample}

\begin{prismbox}[PrismLight]{Öğrenme çıktıları}
Bu bölümün sonunda:
\begin{itemize}
  \item Tip I ve Tip II hatayı araştırma bağlamında ayırt edebilecek,
  \item güç, $\beta$, etki büyüklüğü ve örneklem hacmi arasındaki ilişkiyi açıklayabilecek,
  \item tek örneklem \tstat testini tek ve çift yönlü alternatiflerle uygulayabilecek,
  \item test sonucunu güven aralığı, etki büyüklüğü ve varsayımlarla raporlayabileceksiniz.
\end{itemize}
\end{prismbox}

\begin{prerequisitebox}
$H_0$, $H_1$, anlamlılık düzeyi, \pval-değeri, güven aralığı ve Student $t$
dağılımının temel rolü bilinmelidir.
\end{prerequisitebox}

Hipotez, \pval-değeri ve test yönü kavramları için Bölüm~\ref{ch:hypothesis};
güven aralığı yorumu için Bölüm~\ref{sec:ci-interpretation} temel oluşturur.

\section{Karar ve gerçek durum}
\label{sec:error-types}

\begin{table}[H]
\centering
\caption{Hipotez testinde kararlar ve hata türleri.}
\label{tab:error-decisions}
\begin{tabular}{p{0.25\textwidth}p{0.27\textwidth}p{0.27\textwidth}}
\toprule
 & $H_0$ doğru & $H_0$ yanlış \\
\midrule
$H_0$ reddedilir & Tip I hata ($\alpha$) & Doğru karar (güç) \\
$H_0$ reddedilmez & Doğru karar & Tip II hata ($\beta$) \\
\bottomrule
\end{tabular}
\end{table}

Tablo~\ref{tab:error-decisions}'e göre Tip I hata doğru bir sıfır hipotezini
reddetmek, Tip II hata ise yanlış bir
sıfır hipotezini reddedememektir.

\begin{figure}[H]
\centering
\begin{tikzpicture}[x=1.15cm,y=4.2cm]
  \fill[PrismSoftwarePurple]
    (0,0)--plot[smooth,domain=0:4.4,samples=80]
    (\x,{0.78*exp(-((\x-5)^2)/1.7)})--(4.4,0)--cycle;
  \fill[PrismWarnOrange!80]
    (4.4,0)--plot[smooth,domain=4.4:7.4,samples=80]
    (\x,{0.78*exp(-((\x-3)^2)/1.7)})--(7.4,0)--cycle;
  \draw[PrismBlue,very thick,smooth,domain=0:7.4,samples=100]
    plot (\x,{0.78*exp(-((\x-3)^2)/1.7)});
  \draw[PrismSubheadingBlue,very thick,smooth,domain=0:7.4,samples=100]
    plot (\x,{0.78*exp(-((\x-5)^2)/1.7)});
  \draw[PrismGray,->] (0,0)--(7.6,0);
  \draw[PrismGray,dashed] (4.4,0)--(4.4,0.92)
    node[above,font=\scriptsize]{kritik değer};
  \node[font=\scriptsize,text=PrismBlue] at (2.7,0.84) {$H_0$ doğru};
  \node[font=\scriptsize,text=PrismSubheadingBlue] at (5.5,0.84) {$H_1$ doğru};
  \node[font=\scriptsize] at (5.0,0.18) {Tip I hata $\alpha$};
  \node[font=\scriptsize] at (3.6,0.25) {Tip II hata $\beta$};
\end{tikzpicture}
\caption{Tek yönlü karar kuralında Tip I ve Tip II hata bölgelerinin şematik gösterimi.}
\label{fig:error-regions}
\end{figure}

Bir tarama aracının riskli öğrencileri belirlediği düşünülsün. Gerçekte risk
taşımayan bir öğrenciyi riskli sınıflandırmak Tip I hataya, risk taşıyan bir
öğrenciyi gözden kaçırmak Tip II hataya benzer. İki hatanın sonuçları aynı
olmayabilir; bu nedenle $\alpha$ seçimi yalnız geleneksel bir eşik değil,
yanlış kararların maliyetleriyle ilgili bir araştırma kararıdır.

\section{Test gücü}
\label{sec:power}

Test gücü $1-\beta$, gerçek bir etki varken $H_0$'ı reddetme olasılığıdır.
Etki büyüklüğü ve örneklem hacmi arttıkça, değişkenlik azaldıkça güç genellikle
artar. $\alpha$ düzeyini artırmak gücü yükseltebilir; fakat Tip I hata riskini
de artırır. Güç analizinin etki büyüklüğü, hata oranları ve örneklem hacmiyle
ilişkisi için bkz. \citet{cohen1988} ve \citet{lakens2022}.

\begin{table}[H]
\centering
\caption{Çift yönlü tek örneklem \tstat testinde $d=0{,}50$ için yaklaşık güç.}
\label{tab:power-by-n}
\begin{tabular}{rrrrr}
\toprule
Örneklem hacmi $n$ & 10 & 20 & 34 & 50 \\
\midrule
Güç & 0,29 & 0,56 & 0,81 & 0,93 \\
\bottomrule
\end{tabular}
\end{table}

Tablo~\ref{tab:power-by-n}, orta büyüklükte bir etki için küçük örneklemlerin
gerçek farkı sıklıkla
kaçırabileceğini gösterir. Hesaplar $\alpha=0{,}05$ ve iki yönlü alternatif
içindir; beklenen etki veya tasarım değiştiğinde güç de değişir.

\begin{figure}[H]
\centering
\begin{tikzpicture}[x=0.12cm,y=4.8cm]
  \draw[->,PrismGray] (0,0)--(64,0) node[right,font=\scriptsize]{$n$};
  \draw[->,PrismGray] (0,0)--(0,1.05) node[above,font=\scriptsize]{güç};
  \foreach \x in {10,20,30,40,50,60} {
    \draw[PrismGray] (\x,-0.015)--(\x,0.015);
    \node[below,font=\tiny] at (\x,-0.015) {\x};
  }
  \foreach \y in {0.2,0.4,0.6,0.8,1.0} {
    \draw[PrismGray] (-0.8,\y)--(0.8,\y);
    \node[left,font=\tiny] at (-0.8,\y) {\y};
  }
  \draw[PrismWarnOrange!80!black,dashed] (0,0.8)--(60,0.8)
    node[right,font=\tiny]{0,80};
  \draw[PrismBlue,very thick,smooth]
    plot coordinates {(5,.141)(10,.293)(15,.438)(20,.565)(25,.670)
                      (30,.754)(34,.808)(40,.869)(50,.934)(60,.968)};
  \foreach \x/\y in {5/.141,10/.293,15/.438,20/.565,25/.670,
                       30/.754,34/.808,40/.869,50/.934,60/.968}
    \fill[PrismBlue] (\x,\y) circle (1.6pt);
\end{tikzpicture}
\caption{$d=0{,}50$ için örneklem hacmi arttıkça test gücünün değişimi.}
\label{fig:power-curve}
\end{figure}

Şekil~\ref{fig:power-curve}, aynı hesapları sürekli bir güç eğrisi üzerinde
gösterir ve 0,80 hedefinin yaklaşık $n=34$ çevresinde aşıldığını belirtir.

\section{Tek örneklem \tstat testi}
\label{sec:one-sample-t}

Tek örneklem \tstat testi, bir evren ortalamasını önceden belirlenmiş $\mu_0$
değeriyle karşılaştırır:
\[
t=\frac{\bar x-\mu_0}{s/\sqrt n},\qquad sd=n-1.
\]
Testin temel koşulları gözlemlerin bağımsızlığı ve küçük örneklemlerde ciddi
çarpıklık ya da etkili aykırı değerlerin bulunmamasıdır.

\section{Etki büyüklüğü ve raporlama}

Tek örneklem için standartlaştırılmış fark
\[
d=\frac{\bar x-\mu_0}{s}
\]
ile özetlenebilir. Test sonucu ortalama, standart sapma, t değeri, serbestlik
derecesi, \pval-değeri, güven aralığı ve etki büyüklüğüyle birlikte verilmelidir
\citep{lakens2013}.

\section{Çözümlü örnek}

Bir ölçeğin referans ortalaması 50'dir. $n=25$ öğrenciden elde edilen
$\bar x=55$ ve $s=10$ değerleri için
\[
t=\frac{55-50}{10/\sqrt{25}}=2{,}50,
\qquad sd=24.
\]
Çift yönlü \pval-değeri yaklaşık $0{,}020$'dir. Yüzde 95 güven aralığı
$[50{,}87;59{,}13]$, standartlaştırılmış fark ise $d=0{,}50$'dir. Sonuç şöyle
raporlanabilir: ``Örneklem ortalaması ($\bar x=55$, $s=10$) referans
değerden yüksekti, $t(24)=2{,}50$, $p=0{,}020$, yüzde 95 GA
$[50{,}87;59{,}13]$, $d=0{,}50$.''

Çözümün karar zinciri şöyledir:

\begin{enumerate}
  \item Hedef parametre evren ortalaması, karşılaştırma değeri $\mu_0=50$'dir.
  \item Gözlemler bağımsız kabul edilir; puan dağılımında etkili aykırı değer bulunmadığı doğrulanır.
  \item Standart hata $10/\sqrt{25}=2$ ve test istatistiği $2{,}50$ hesaplanır.
  \item $sd=24$ için \pval-değeri $0{,}020$ olduğundan $H_0$ reddedilir.
  \item Güven aralığı farkın 0,87 ile 9,13 puan arasında olabileceğini gösterir.
  \item $d=0{,}50$ sonucu bağlama ve ölçme aracının birimine göre değerlendirilir.
\end{enumerate}

\begin{assumptionbox}
Tek örneklem \tstat testi için referans değerin bilimsel olarak anlamlı ve bağımsız
bir kaynaktan belirlenmiş olması gerekir. Örneklem ortalamasına bakıldıktan
sonra uygun görünen bir $\mu_0$ seçmek testin hata oranını bozar.
\end{assumptionbox}

\begin{mistakebox}
``Anlamlı sonuç elde etmek'' amacıyla veri toplama sürecini \pval-değeri
0,05'in altına düşünce durdurmak Tip I hata oranını artırır. Örneklem hacmi
ve durdurma kuralı analizden önce belirlenmelidir.
\end{mistakebox}

\section{Güç için planlama mantığı}

Güç analizi beklenen etki büyüklüğü, seçilen $\alpha$, hedef güç ve test
yönüne dayanır. Çok iyimser etki varsayımı yetersiz örnekleme yol açabilir.
Bu nedenle önceki çalışmalar, pilot veri ve bilimsel olarak anlamlı en küçük
etki birlikte değerlendirilmelidir.

\section{Python ile yeniden üretim}

\begin{softwarebox}
\begin{lstlisting}
from scipy import stats
import numpy as np

n, mean, sd, mu0 = 25, 55, 10, 50
se = sd / np.sqrt(n)
t_value = (mean - mu0) / se
p_value = 2 * stats.t.sf(abs(t_value), df=n-1)
critical = stats.t.ppf(.975, df=n-1)
ci = (mean - critical * se, mean + critical * se)

print(t_value, p_value, ci)

alpha, target_n, d = .05, 34, .50
df = target_n - 1
critical = stats.t.ppf(1 - alpha / 2, df=df)
ncp = d * np.sqrt(target_n)
power = (stats.nct.cdf(-critical, df=df, nc=ncp)
         + stats.nct.sf(critical, df=df, nc=ncp))
print("Guc:", power)
\end{lstlisting}
Kod çıktısı yaklaşık olarak $t=2{,}50$, $p=0{,}0197$ ve
$[50{,}87;59{,}13]$ güven aralığını verir. İkinci hesap $d=0{,}50$ için
$n=34$ durumunda gücün yaklaşık 0,808 olduğunu doğrular. Yuvarlama yalnız
raporlama aşamasında yapılmalıdır.
\end{softwarebox}

\section{Çözümlü problemler}

\subsection*{Problem 1: Tip I ve Tip II hatanın bağlamsal anlamı}

Bir üniversite, yeni bir akademik destek programının öğrencilerin ortalama
başarı puanını değiştirmediği $H_0$ hipotezini sınamaktadır. Bu araştırmada
Tip I ve Tip II hatayı açıklayın; olası sonuçlarını karşılaştırın.

\textbf{Çözüm.} Tip I hata, programın gerçekte ortalama başarıyı
değiştirmediği hâlde test sonucunda etkili olduğuna karar vermektir. Üniversite
bu hata sonucunda etkisiz bir programa bütçe, zaman ve personel ayırabilir.

Tip II hata, programın gerçekte başarı üzerinde etkisi olduğu hâlde testin
bu etkiyi belirleyememesi ve $H_0$'ın reddedilememesidir. Bu durumda yararlı
bir program uygulanmadan bırakılabilir. İki hatanın maliyeti bağlama göre
farklıdır; bu nedenle $\alpha$, hedef güç ve bilimsel olarak önemli en küçük
etki çalışma başlamadan önce birlikte değerlendirilmelidir.

\textbf{Sık hata.} Tip I hata ``hesaplama hatası'', Tip II hata ise ``veri
giriş hatası'' değildir. İkisi de belirsizlik altında verilen test kararının
gerçek durumla uyuşmamasını ifade eder.

\subsection*{Problem 2: Çift yönlü tek örneklem \tstat testi}

Bir ölçeğin referans ortalaması 38'dir. Bağımsız 16 öğrencinin ortalaması 42,
standart sapması 8 bulunmuştur. Ortalama puanın 38'den farklı olup olmadığını
yüzde 5 düzeyinde sınayın; yüzde 95 güven aralığını ve Cohen'in $d$ değerini
hesaplayın.

\textbf{Çözüm.} Hipotezler
\[
H_0:\mu=38,
\qquad H_1:\mu\neq38
\]
biçimindedir. Standart hata ve test istatistiği
\[
SE=\frac{8}{\sqrt{16}}=2,
\qquad
t=\frac{42-38}{2}=2{,}00,
\qquad sd=15
\]
olur. Çift yönlü \pval-değeri yaklaşık 0,064'tür; dolayısıyla $H_0$ yüzde 5
düzeyinde reddedilemez.

$t_{0{,}975,15}=2{,}131$ kullanıldığında ortalama için güven aralığı
\[
42\pm2{,}131(2)=[37{,}74;46{,}26]
\]
olur ve referans değer 38'i içerir. Standartlaştırılmış fark
\[
d=\frac{42-38}{8}=0{,}50
\]
olarak bulunur. Orta büyüklükteki gözlenen etkiye rağmen küçük örneklem
belirsizliği geniştir.

\textbf{Raporlama örneği.} ``Örneklem ortalaması 42 ($s=8$) olmakla birlikte
38 referans değerinden fark yüzde 5 düzeyinde belirlenemedi,
$t(15)=2{,}00$, $p=0{,}064$, yüzde 95 GA $[37{,}74;46{,}26]$,
$d=0{,}50$.''

\subsection*{Problem 3: Aynı veride tek yönlü test}

Önceki problemde araştırma hipotezinin veri toplanmadan önce
$H_1:\mu>38$ olarak belirlendiğini varsayın. Kararı yeniden değerlendirin.

\textbf{Çözüm.} Test istatistiği ve serbestlik derecesi değişmez:
$t(15)=2{,}00$. Üst yönlü \pval-değeri çift yönlü değerin yarısıdır ve yaklaşık
0,032'dir. Bu nedenle önceden belirlenmiş $\alpha=0{,}05$ düzeyinde $H_0$
reddedilir; ortalamanın 38'in üzerinde olduğuna ilişkin kanıt vardır.

Bu sonuç, veri görüldükten sonra tek yönlü test seçilebileceği anlamına
gelmez. Azalma yönündeki bilimsel olarak önemli bir sonuç da araştırma
sorusunun parçasıysa çift yönlü test korunmalıdır. Tek yönlü kararın güven
aralığı da yönle uyumlu tek taraflı sınır üzerinden raporlanabilir.

\subsection*{Problem 4: Güç ve beklenen Tip II hata sayısı}

$d=0{,}50$, $\alpha=0{,}05$ ve çift yönlü tek örneklem \tstat testi için $n=34$
durumunda güç 0,808 olsun. $\beta$ değerini bulun ve aynı gerçek etki altında
1.000 bağımsız araştırmada beklenen karar sayılarını yorumlayın.

\textbf{Çözüm.} Tip II hata olasılığı
\[
\beta=1-\text{güç}=1-0{,}808=0{,}192
\]
olur. Aynı koşullarda 1.000 bağımsız çalışma yürütülse uzun dönemde yaklaşık
808 çalışmanın $H_0$'ı reddederek etkiyi belirlemesi, yaklaşık 192 çalışmanın
ise gerçek etkiyi kaçırması beklenir.

Bu sayılar belirli 1.000 çalışmada kesin olarak gerçekleşecek frekanslar
değildir; uzun dönemli beklentilerdir. Yüzde 80 güç, tek bir çalışmanın doğru
sonuç vereceğini garanti etmez.

\subsection*{Problem 5: Yaklaşık güç temelli örneklem planı}

Çift yönlü yüzde 5 anlamlılık düzeyinde, standartlaştırılmış $d=0{,}50$
etkisini yüzde 80 güçle belirlemek için normal yaklaşım kullanılsın.
$z_{0{,}975}=1{,}96$ ve $z_{0{,}80}=0{,}842$ alarak yaklaşık örneklem hacmini
hesaplayın.

\textbf{Çözüm.} Tek grup ortalaması için yaklaşık formül
\[
n\approx\left(\frac{z_{1-\alpha/2}+z_{1-\beta}}{d}\right)^2
\]
biçimindedir. Değerler yerine yazılırsa
\[
n\approx\left(\frac{1{,}96+0{,}842}{0{,}50}\right)^2
=31{,}40
\]
bulunur ve yukarı yuvarlanarak 32 kişi elde edilir. Ancak bu normal yaklaşım
$\sigma$ biliniyormuş gibi davranır. Standart sapmanın örneklemden tahmin
edildiği kesin t hesabı biraz daha büyük bir örneklem gerektirir; bölümdeki
güç tablosunda yaklaşık yüzde 80 güç için $n=34$ görülmektedir.

Beklenen kayıp veya dışlama varsa başlangıç örneklem hacmi ayrıca
artırılmalıdır. Güç planı gözlenen etki üzerinden analiz sonrasında
gerekçelendirilmemelidir.

\subsection*{Problem 6: Büyük örneklemde küçük etki}

Referans ortalama 50 ve örneklem standart sapması 10 olsun. Hem $n=100$ hem
$n=10\,000$ için örneklem ortalaması 51,5 bulunmuştur. Test istatistiklerini,
yaklaşık sonuçları ve etki büyüklüğünü karşılaştırın.

\textbf{Çözüm.} Her iki çalışmada standartlaştırılmış fark aynıdır:
\[
d=\frac{51{,}5-50}{10}=0{,}15.
\]
$n=100$ için $SE=10/\sqrt{100}=1$ ve
\[
t=\frac{1{,}5}{1}=1{,}50.
\]
Çift yönlü \pval-değeri yaklaşık 0,137'dir; $H_0$ reddedilemez. Yaklaşık yüzde 95
güven aralığı fark için $1{,}5\pm1{,}984(1)$, yani
$[-0{,}48;3{,}48]$'dir.

$n=10\,000$ için $SE=10/100=0{,}10$ ve
\[
t=\frac{1{,}5}{0{,}10}=15
\]
olur; \pval-değeri çok küçüktür. Yaklaşık yüzde 95 fark aralığı
$1{,}5\pm1{,}96(0{,}10)=[1{,}30;1{,}70]$'tir.

Büyük örneklem aynı küçük etkiyi çok hassas biçimde sıfırdan ayırmıştır;
ancak etkinin büyüklüğü iki çalışmada da $d=0{,}15$'tir. İstatistiksel kanıt
güçlenirken uygulamalı önem kendiliğinden büyümemiştir.

\begin{outputguidebox}
Tek örneklem \tstat testi çıktısında örneklem ortalaması ve standart sapmasını,
referans değeri, ortalama farkını, t değerini, serbestlik derecesini,
\pval-değerini ve güven aralığını birlikte okuyun. Güç analizinde kullanılan
etki, yön, $\alpha$ ve hedef güç değerlerinin çalışma öncesinde belirlendiğini
doğrulayın.
\end{outputguidebox}

\section{Literatür verisiyle yazılım uygulamaları}
\label{sec:spss-b10}

\textbf{Amaç ve veri.} \texttt{b10.csv} ile anlamlılık, etki büyüklüğü
ve ileriye dönük örneklem planını ayırın \citep{cortez2008data}.
Gözlenen analizde $H_0:\mu=10$ ve $H_1:\mu\ne10$ sınanır.
Planlamada ise yeni bir çalışma için hedef fark 1 puandır;
bu fark, gözlenen ortalama farkından ayrı bir girdidir.

\textbf{Ortak veri.} Doğrulanan B10 uygulamasında üç yazılım da
\nolinkurl{companion/spss/csv/b10.csv} dosyasını kullanmıştır.
Dosya 395 satır ve 10 sütun içerir; kayıtların tamamı analiz edilir.
B06'nın seçim göstergeleri ve ağırlığı burada kullanılmaz.
Bu karşılaştırma Excel içe aktarmanın doğrulandığı anlamına gelmez.
Kodları \texttt{companion} klasörünü içeren paket kökünden çalıştırın;
veri sözlüğü ve hazırlık için bkz. sayfa~\pageref{sec:spss-guide}.

\subsection{IBM SPSS uygulaması}
\textbf{Başlamadan önce.} Aşağıdaki kod CSV verisini yeniden açar.
Menü yolunu kullanacaksanız veri açma ve tanımlama komutlarını
\texttt{EXECUTE} satırı dahil çalıştırın. Filtre, ağırlık ve dosya
bölme kapalı olmalıdır. \texttt{AGGREGATE}, etkin veriyi tek özet
satırıyla değiştirir; bu özeti ham veri dosyasının üzerine kaydetmeyin.
Tekrar analiz için veri açma adımından başlayın. Güç komutu SPSS
Statistics 27 ve sonrasında bulunur; eski sürümlerde yalnız bu adımı
atlayıp R veya Python planını kullanabilirsiniz.

\begin{spssadimlar}
\spssadim{Gözlenen test ve $d$ için aşağıdaki kodu kullanın. \emph{One-Sample Test} tablosunda \emph{Two-Sided p} sütununu, \texttt{LIST} çıktısında \texttt{etki\_d} değerini okuyun. Örneklem planını bu gözlenen farka göre kurmayın.}
\spssadim{\emph{Analyze $\to$ Power Analysis $\to$ Means $\to$ One-Sample T Test} yolunu açın. \emph{Estimate} için \emph{Sample size} seçin.}
\spssadim{Hedef gücü \texttt{0.80} girin. \emph{Hypothesized Values} ile \emph{Population mean}=11 ve \emph{Null value}=10 belirleyin.}
\spssadim{\emph{Population standard deviation} alanına \texttt{4.581442611} girin; \emph{Nondirectional (two-sided)} ve anlamlılık düzeyi \texttt{0.05} seçin.}
\spssadim{\emph{OK} ile gereken örneklem hacmini alın. Paylaşılan çıktıda \emph{N}=167 ve \emph{Actual Power}=0,801 görülür. Bu güç, planlanan 1 puan fark içindir; eldeki 395 kaydın gözlenen gücü değildir.}
\end{spssadimlar}

{\raggedright
\noindent\textbf{Komutların görevi.} \texttt{T-TEST} ve \texttt{COMPUTE etki\_d} gözlenen veriyi inceler. \texttt{POWER MEANS ONESAMPLE} ayrı bir ileriye dönük planlamadır; hedef güç ve alternatif ortalama önceden belirlenir.\par}

\textbf{Uygulama kodu:} \texttt{b10}. Doğrulanan veri dosyası
\texttt{companion/spss} altındaki \texttt{csv/b10.csv} dosyasıdır.

\begin{lstlisting}[style=prismSPSS]
SET DECIMAL=DOT.
GET DATA /TYPE=TXT
 /FILE='companion/spss/csv/b10.csv'
 /ENCODING='UTF8'
 /ARRANGEMENT=DELIMITED /DELCASE=LINE
 /DELIMITERS="," /QUALIFIER='"'
 /FIRSTCASE=2
 /VARIABLES=id F8.0 school F8.0 sex F8.0 age F8.0
 studytime F8.0 absences F8.0 g1 F8.0 g2 F8.0
 g3 F8.0 pass10 F8.0.
FILTER OFF.
WEIGHT OFF.
SPLIT FILE OFF.
VARIABLE LABELS g3 'Yil sonu matematik notu'.
VARIABLE LEVEL
 id school sex pass10 (NOMINAL)
 studytime (ORDINAL)
 age absences g1 g2 g3 (SCALE).
FORMATS g3 (F14.9).
EXECUTE.
\end{lstlisting}

\par\noindent\begin{minipage}{\linewidth}
\textbf{Gözlenen test ve özet.} Veri açma bloğundan sonra çalıştırın.

\begin{lstlisting}[style=prismSPSS]
T-TEST /TESTVAL=10 /VARIABLES=g3
 /CRITERIA=CI(.95).
AGGREGATE /OUTFILE=*
 /BREAK= /n=N /ort=MEAN(g3) /ss=SD(g3).
COMPUTE etki_d=(ort-10)/ss.
FORMATS n (F8.0).
FORMATS ort ss etki_d (F14.9).
LIST VARIABLES=n ort ss etki_d.
\end{lstlisting}
\end{minipage}\par

\par\noindent\begin{minipage}{\linewidth}
\textbf{Ayrı planlama adımı.} Aşağıdaki komut ham veriden yeni bir
etki tahmin etmez; belirtilen varsayımlarla gereken hacmi hesaplar.

\begin{lstlisting}[style=prismSPSS]
POWER MEANS ONESAMPLE
 /PARAMETERS TEST=NONDIRECTIONAL
 SIGNIFICANCE=.05 POWER=.80
 SD=4.581442611 MEAN=11 NULL=10.
\end{lstlisting}

\end{minipage}\par

\begin{outputguidebox}
\textbf{Paylaşılan SPSS çıktısıyla doğrulanan değerler.}
\emph{One-Sample Statistics} tablosunda $n=395$, ortalama
10,415189873, standart sapma 4,5814426110 ve standart hata
0,23051739492 görülür. Test tablosunda $t(394)=1{,}801$ ve
iki yönlü $p=0{,}072$; \texttt{LIST} çıktısında
$d=(\bar x-10)/s=0{,}090624266$ bulunmuştur.
Tek yönlü 0,036 değeri bu uygulamanın kararında kullanılmaz.

Güç tablosundaki \emph{Effect Size}=0,218 ise
$1/4{,}581442611$ planlama etkisinin yuvarlanmış değeridir;
gözlenen $d$ ile aynı değildir. \emph{Power}=0,8 hedef gücü,
\emph{Actual Power}=0,801 ise 167 kişide hesaplanan gücü gösterir.
\texttt{LIST} altındaki bir kayıt, 395 kişinin toplu özetidir;
örneklem büyüklüğünün bire düştüğü anlamına gelmez.
\end{outputguidebox}


\par\noindent\begin{minipage}{\linewidth}
\subsection{R uygulaması}
\label{sec:r-gercek-b10}

\textbf{Gözlenen analiz.} Standartlaştırılmış fark ile iki yönlü
$p$ değeri ayrı niceliklerdir. Ortalama aralığından 10 çıkarılarak
SPSS'in fark aralığı elde edilir.

\begin{lstlisting}[style=prismR]
veri <- read.csv("companion/spss/csv/b10.csv")
stopifnot(nrow(veri) == 395, ncol(veri) == 10,
  !anyNA(veri),
  all(veri$pass10 == as.integer(veri$g3 >= 10)))
sonuc <- t.test(veri$g3, mu=10,
  alternative="two.sided", conf.level=.95)
ortalama <- mean(veri$g3)
standart_sapma <- sd(veri$g3)
print(c(n=nrow(veri), ort=ortalama, ss=standart_sapma,
  sh_ort=standart_sapma/sqrt(nrow(veri)),
  t=unname(sonuc$statistic), df=unname(sonuc$parameter),
  p_iki_yonlu=sonuc$p.value, fark=ortalama-10,
  etki_d=(ortalama-10)/standart_sapma,
  ort_alt=sonuc$conf.int[1],
  ort_ust=sonuc$conf.int[2]), digits=15)
print(sum(veri$g3 == 0))
\end{lstlisting}

\textbf{Çıktıyı okuma.} Paylaşılan R çıktısında 395 satır, 10 sütun,
sıfır eksik değer ve analizde tutulan 38 sıfır not vardır.
Gözlenen etki 0,090624, iki yönlü $p$ değeri 0,072448'dir.
\end{minipage}\par

\par\noindent\begin{minipage}{\linewidth}
\textbf{İleriye dönük plan.} \texttt{strict=TRUE} iki kuyruğu da
hesaba katar; \texttt{ceiling} gereken hacmi yukarı yuvarlar.

\begin{lstlisting}[style=prismR]
plan <- power.t.test(delta=1, sd=4.581442611,
  sig.level=.05, power=.80, type="one.sample",
  alternative="two.sided", strict=TRUE)
n_plan <- ceiling(plan$n)
guc <- function(hacim) {
  power.t.test(n=hacim, delta=1, sd=4.581442611,
    sig.level=.05, type="one.sample",
    alternative="two.sided", strict=TRUE)$power
}
print(c(n_kesirli=plan$n, n_plan=n_plan,
  guc_plan=guc(n_plan), n_bir_eksik=n_plan-1,
  guc_bir_eksik=guc(n_plan-1)), digits=15)
\end{lstlisting}

\textbf{Çıktıyı okuma.} R çıktısında kesirli hacim yaklaşık
166,6755, yukarı yuvarlanan hacim 167'dir. Güç 166 kişide
0,798386, 167 kişide 0,800771 olur; hedefe aşağı yuvarlanmaz.
\end{minipage}\par

\par\noindent\begin{minipage}{\linewidth}
\subsection{Python uygulaması}
\label{sec:python-b10}

\textbf{Gözlenen analiz.} \texttt{ddof=1} örneklem standart
sapmasını seçer. \texttt{aralik} ortalamanın yüzde 95 aralığıdır.

\begin{lstlisting}[style=prismPythonApp]
import pandas as pd
import numpy as np
from scipy import stats, optimize

veri = pd.read_csv("companion/spss/csv/b10.csv")
assert veri.shape == (395, 10)
assert not veri.isna().any().any()
assert veri.pass10.eq((veri.g3 >= 10).astype(int)).all()
sonuc = stats.ttest_1samp(veri.g3, 10,
                         alternative="two-sided")
aralik = sonuc.confidence_interval(.95)
ortalama = veri.g3.mean()
standart_sapma = veri.g3.std(ddof=1)
ozet = dict(n=len(veri), ort=ortalama, ss=standart_sapma,
    sh_ort=standart_sapma/np.sqrt(len(veri)),
    t=sonuc.statistic, df=len(veri)-1,
    p_iki_yonlu=sonuc.pvalue, fark=ortalama-10,
    etki_d=(ortalama-10)/standart_sapma,
    ort_alt=aralik.low, ort_ust=aralik.high)
for ad, deger in ozet.items():
    print(f"{ad}: {deger:.13f}")
print(int((veri.g3 == 0).sum()))
\end{lstlisting}

\textbf{Çıktıyı okuma.} Paylaşılan Python çıktısı da 395 satır,
10 sütun, sıfır eksik değer ve 38 sıfır not gösterir. Gözlenen
analizin 11 özet değeri R ile 13 ondalık basamakta eşleşir.
R ve Python kodları ham veriyi tek satırlık özetle değiştirmez.
\end{minipage}\par

\par\noindent\begin{minipage}{\linewidth}
\textbf{İleriye dönük plan.} Önceki Python bloğundan sonra çalıştırın.
\texttt{nct}, merkezî olmayan $t$ dağılımıyla iki kuyruklu gücü
hesaplar. Kök çözümü yukarı yuvarlanır; bir eksik kişi de kontrol edilir.

\begin{lstlisting}[style=prismPythonApp]
def guc(hacim):
    serbestlik = hacim - 1
    kritik = stats.t.ppf(.975, serbestlik)
    kayma = np.sqrt(hacim) / 4.581442611
    return (stats.nct.cdf(-kritik, serbestlik, kayma)
            + stats.nct.sf(kritik, serbestlik, kayma))

n_kesirli = optimize.brentq(
    lambda hacim: guc(hacim)-.80, 2, 1000)
n_plan = int(np.ceil(n_kesirli))
assert guc(n_plan-1) < .80 <= guc(n_plan)
print("n_kesirli:", n_kesirli, "n_plan:", n_plan)
print("guc_plan:", guc(n_plan))
print("n_bir_eksik:", n_plan-1,
      "guc_bir_eksik:", guc(n_plan-1))
\end{lstlisting}

\textbf{Çıktıyı okuma.} Python da 167 kişi sonucunu verir.
Planlama standart sapması üç yazılımda da tam olarak
\texttt{4.581442611} girilmiştir; hesap öncesinde 4,58'e yuvarlanmaz.
\end{minipage}\par

\subsection{Sonuçların karşılaştırılması ve ortak rapor}
13 Eylül 2026'da paylaşılan SPSS tabloları ile R ve Python metin
çıktıları karşılaştırılmıştır. Gözlenen analizin 11 özet değeri
R ve Python'da 13 ondalık basamakta eşleşir; SPSS değerleri gösterilen
yuvarlama hassasiyetlerinde uyumludur. R ve Python'un kesirli hacimleri
arasında yaklaşık $1{,}1\times10^{-7}$, güçlerinde yaklaşık
$2\times10^{-12}$ düzeyinde sayısal farklar vardır. Bu farklar
167 kişilik planı değiştirmez; bütün planlama sonuçları için
13 basamakta eşitlik iddia edilmez. Proje CSV'sinden yapılan hesaplar
da sonuçları destekler. Aşağıdaki tablolar çıktıların kitap düzenine
uyarlanmış özetidir; ekran görüntüsü değildir.

\begin{table}[H]
\centering
\caption{B10'da gözlenen test ve standartlaştırılmış fark}
\label{tab:b10-dogrulanmis-gozlenen}
\begin{tabular}{@{}lr@{}}
\toprule
Ölçü & Değer \\
\midrule
Gözlem sayısı & 395 \\
Ortalama not & 10,415190 \\
Notların standart sapması & 4,581443 \\
Ortalamanın standart hatası & 0,230517 \\
Referanstan fark ($\bar x-10$) & 0,415190 \\
Gözlenen etki ($d$) & 0,090624 \\
$t$ istatistiği & 1,801122 \\
Serbestlik derecesi & 394 \\
İki yönlü $p$ değeri & 0,072448 \\
Ortalama için \%95 güven aralığı & $[9{,}961992;\ 10{,}868388]$ \\
Fark için \%95 güven aralığı & $[-0{,}038008;\ 0{,}868388]$ \\
\bottomrule
\end{tabular}
\par\smallskip
{\footnotesize Ondalıklı değerler altı basamağa yuvarlanmıştır.
Standart sapmada 394 böleni kullanılmıştır. Fark aralığı, ortalama
aralığının her iki sınırından 10 çıkarılarak elde edilir.\par}
\end{table}

\begin{table}[H]
\centering
\caption{B10'da bir puanlık hedef fark için örneklem planı}
\label{tab:b10-dogrulanmis-plan}
\begin{tabular}{@{}lr@{}}
\toprule
Planlama girdisi veya sonucu & Değer \\
\midrule
Sıfır hipotezindeki ortalama & 10 \\
Planlanan alternatif ortalama & 11 \\
Hedef fark (puan) & 1 \\
Planlama standart sapması & 4,581442611 \\
Planlama etkisi ($1/4{,}581442611$) & 0,218272 \\
İki yönlü anlamlılık düzeyi & 0,05 \\
Hedef güç & 0,80 \\
Kesirli örneklem hacmi (yaklaşık) & 166,6755 \\
Gereken en küçük tam sayı hacim & 167 \\
166 kişide hesaplanan güç & 0,798386 \\
167 kişide hesaplanan güç & 0,800771 \\
\bottomrule
\end{tabular}
\par\smallskip
{\footnotesize Kesirli hacim R ve Python hesaplarından gelir.
SPSS doğrudan 167 kişi ve üç basamakla 0,801 güç göstermiştir.
Güç hesabı bağımsız normal gözlemler modelinde merkezî olmayan
$t$ dağılımının iki kuyruğunu kullanır.\par}
\end{table}

\Needspace{13\baselineskip}
\begin{outputguidebox}
Tablo~\ref{tab:b10-dogrulanmis-gozlenen}, eldeki 395 kaydı inceler.
İki yönlü $p=0{,}072448>0{,}05$ olduğundan $H_0$ reddedilmez;
bu, ortalamanın kesinlikle 10 olduğunun veya farkın eğitim bakımından
önemsiz olduğunun kanıtı değildir. Sıfır notlu 38 kayıt korunmuştur.

Tablo~\ref{tab:b10-dogrulanmis-plan} ise farklı bir soruya yanıt
verir: 1 puanlık gerçek farkı, belirtilen standart sapma ve model
altında yüzde 80 güçle saptamak için kaç kişi gerekir?
Gözlenen $d=0{,}090624$ yerine planlama etkisi 0,218272 kullanılır.
167 kişideki 0,800771 güç, gözlenen $p$ değerinden türetilmiş
geriye dönük güç değildir; 395 kaydın 167'sini seçme önerisi de değildir.

Bir puanın anlamlı hedef sayılması ve planlama standart sapmasının
yeni çalışmaya uygunluğu ayrıca gerekçelendirilmelidir. Kayıplar ve
kümelenme bu hesaba dahil değildir. İki okuldan gelen veride yazılımların
uyumu, bağımsızlığı veya daha geniş bir evrene genellenebilirliği kanıtlamaz.
\end{outputguidebox}

\Needspace{10\baselineskip}
\textbf{Raporlama örneği (doğrulanmış çıktılarla).}
``395 kayıtta ortalama not 10,415190 ve standart sapma 4,581443
bulunmuştur. 10 puan referansına karşı iki yönlü testte
$t(394)=1{,}801122$, $p=0{,}072448$ ve standartlaştırılmış fark
$d=0{,}090624$ elde edilmiştir. Sıfır hipotezi yüzde 5 düzeyinde
reddedilmemiştir. Ayrı bir ileriye dönük planlamada, 1 puanlık hedef
fark, 4,581442611 standart sapma, iki yönlü 0,05 anlamlılık düzeyi
ve 0,80 hedef güç için en az 167 kişi gerektiği hesaplanmıştır;
bu hacimde güç 0,800771'dir. Plan bağımsız normal gözlemler modeline
dayanır; kayıp ve kümelenme etkilerini içermez. İki okulla sınırlı
kaynak verinin temsil gücü bu hesaplarla doğrulanmış değildir.''


\textbf{Görev.} Bir puanlık farkın eğitim bakımından neden önemli kabul
edilebileceğini gerekçelendirin. Bu gerekçe yoksa yalnız güç hesabı yapmak
çalışmanın bilimsel önemini sağlamaz. Gözlenen $p$ değerinden geriye dönük
güç üretip sonucu yeniden değerlendirmeyin. 166 kişiye aşağı yuvarlamanın
neden hedefi karşılamadığını açıklayın. SPSS'te sonraki uygulamaya
geçerken tek satırlık özet yerine gerekli ham veriyi yeniden açın.

\section{R ile test gücü ve tek örneklem testi}
\label{sec:r-b10}

\textbf{Soru.} $n=25$, $\bar x=55$, $s=10$ için referans ortalama 50'yi
çift yönlü test edin. Ayrıca veri toplanmadan önce gerçek farkın 5, evren
standart sapmasının 10 olduğu varsayımıyla $n=34$ için gücü hesaplayın.

\begin{lstlisting}[style=prismR]
hacim <- 25
ortalama <- 55
standart_sapma <- 10
referans <- 50
standart_hata <- standart_sapma / sqrt(hacim)
t_degeri <- (ortalama - referans) / standart_hata
p_degeri <- 2 * pt(abs(t_degeri), df = hacim - 1,
                   lower.tail = FALSE)
aralik <- ortalama + c(-1, 1) *
  qt(0.975, df = hacim - 1) * standart_hata
c(t = t_degeri, p = p_degeri,
  d = (ortalama - referans) / standart_sapma)
aralik
power.t.test(n = 34, delta = 5, sd = 10,
             sig.level = 0.05, type = "one.sample",
             alternative = "two.sided", strict = TRUE)
plan <- power.t.test(delta = 5, sd = 10, power = 0.80,
                     sig.level = 0.05, type = "one.sample",
                     alternative = "two.sided", strict = TRUE)
ceiling(plan$n)
\end{lstlisting}

\begin{outputguidebox}
\textbf{Çözüm ve kontrol değerleri.} $t(24)=2{,}50$, $p=0{,}01965$,
$d=0{,}50$ ve yüzde 95 güven aralığı $[50{,}8722;59{,}1278]$'dir.
Planlama varsayımları altında 34 gözlemin gücü yaklaşık 0,8078'dir.
Yüzde 80 hedef güç için hesaplanan hacim yukarı yuvarlandığında 34 bulunur.
\end{outputguidebox}

\textbf{Yorum.} \texttt{strict = TRUE} çift yönlü testin iki ret bölgesini
de güç hesabına katar. Güç hesabındaki 5 puan ve 10 puan önceden gerekçelenmiş
planlama değerleri olmalıdır; gözlenen anlamlılığı yeniden ifade eden
``gözlenen güç'' olarak sunulmamalıdır.

\textbf{Pekiştirme ve yanıt.} Tip II hata olasılığı bu belirli alternatif
ve tasarım altında yaklaşık $1-0{,}8078=0{,}1922$'dir. Daha küçük bir gerçek
fark için aynı örneklem hacminin gücü daha düşük olur.

\section{Bölüm özeti}

\begin{itemize}
  \item Tip I ve Tip II hata birbirinden farklı karar riskleridir.
  \item Güç, gerçek bir etkiyi belirleme olasılığıdır ve çalışma öncesinde planlanmalıdır.
  \item Tek örneklem \tstat testi bir ortalamayı önceden belirlenmiş referans değerle karşılaştırır.
  \item Test sonucu \pval-değeri, güven aralığı ve etki büyüklüğüyle birlikte raporlanır.
  \item Tek ve çift yönlü test seçimi veri görülmeden önce araştırma sorusuna
  göre yapılmalıdır.
  \item Büyük örneklem küçük etkileri anlamlı kılabilir; etki büyüklüğü ve
  uygulamalı önem ayrıca değerlendirilmelidir.
\end{itemize}

\section*{Alıştırmalar}

\begin{enumerate}
  \item Tip I ve Tip II hataya birer örnek verin.
  \item Test gücünü etkileyen üç etken yazın.
  \item Tek örneklem \tstat testinin araştırma sorusunu açıklayın.
  \item Çözümlü örnekte $n=100$ olsaydı standart hata ve t değeri nasıl değişirdi?
  \item Güç analizinde yalnız \pval-değerine odaklanmanın neden yetersiz olduğunu açıklayın.
  \item Güç tablosunda $n=10$ ile $n=50$ sonuçlarını Tip II hata bakımından karşılaştırın.
  \item $n=16$, $\bar x=42$, $s=8$ ve $\mu_0=38$ için t değerini hesaplayın.
  \item Referans değerin veri görüldükten sonra seçilmesinin neden sorunlu olduğunu açıklayın.
  \item Bir tarama testinde yanlış alarm ve gerçek riski gözden kaçırma
  durumlarını Tip I ve Tip II hatayla eşleştirin.
  \item Güç 0,90 ise $\beta$ değerini ve 500 gerçek etkili çalışmada beklenen
  Tip II hata sayısını bulun.
  \item $n=36$, $\bar x=74$, $s=12$ ve $\mu_0=70$ için t değerini ve
  Cohen'in $d$ değerini hesaplayın.
  \item Önceki sorudaki test çift yönlü ve üst yönlü olduğunda \pval-değerlerinin
  nasıl ilişkileneceğini açıklayın.
  \item $d=0{,}40$, $\alpha=0{,}05$ ve hedef güç 0,80 için normal yaklaşım
  formülüyle yaklaşık örneklem hacmini hesaplayın.
  \item Etki büyüklüğü küçüldükçe aynı gücü korumak için örneklem hacminin
  neden artması gerektiğini açıklayın.
  \item Aynı ortalama fark ve standart sapma için $n=25$ ile $n=400$
  çalışmalarının t değerlerini ve \pval-değerlerini nitel olarak karşılaştırın.
  \item İstatistiksel olarak anlamlı fakat uygulamada önemsiz bir tek örneklem
  \tstat testi sonucu tasarlayıp uygun raporlama cümlesi yazın.
\end{enumerate}